\section{部分積分の公式}\label{節:部分積分の導出}

一般的な形は
\[
 \int u(x)v'(x)\dd x
\]
で与えられる.ここで $u$ は微分すると簡単になる因子,\,$v'$ は原始関数 $v$ が分かる因子である.この積分を,\,$u$ を微分し $v'$ を積分した積 $u'v$ の積分に帰着させることを目標とする.

2つの関数 $u(x),v(x)$ が微分できるとき,積の微分は
\[
 \{u(x)v(x)\}'=u'(x)v(x)+u(x)v'(x)
\]
である.求めたいのは右辺の第2項の積分なので,第1項を移項して
\[
 u(x)v'(x)=\{u(x)v(x)\}'-u'(x)v(x)
\]
と書き直してやる.両辺の原始関数を考えると,右辺の第1項からは積 $uv$ が,第2項からはもう1つの積分が現れる.

\begin{theorem}{部分積分の公式}
$u,v$ がある区間で連続微分可能ならば,その区間で
\[
 \int u(x)v'(x)\dd x
 =u(x)v(x)-\int u'(x)v(x)\dd x
\]
が成り立つ.不定積分は任意定数を含めた原始関数の族として読む.特に,同じ区間の $a,b$ に対して
\[
 \int_a^b u(x)v'(x)\dd x
 =[u(x)v(x)]_a^b-\int_a^b u'(x)v(x)\dd x
\]
が成り立つ.
\end{theorem}

定積分の公式は,積の微分を $a$ から $b$ まで積分し,基本定理を使ったものである.\,$[uv]_a^b$ は $u(b)v(b)-u(a)v(a)$ を表す.積の微分のうち一方の項だけを移項したので,この端点の項を省くことはできない.

公式の左辺では $u$ をそのまま,\,$v'$ を微分された形で書いている.右辺へ移ると,\,$u$ は $u'$ へ微分され,\,$v'$ は原始関数 $v$ に変わる.したがって実際に考えることは,「どの因子を微分し,どの因子を積分すると,残る積 $u'v$ が扱いやすくなるか」である.

ときには $\dd v=v'(x)\dd x$,\,$\dd u=u'(x)\dd x$ とまとめて,
\[
 \int u\dd v=uv-\int v\dd u
\]
とも書く.これは同じ公式の短い表記であり,\,$u$ と $v$ を関係のない変数として自由に動かしているのではない.本章では,微分する因子と積分する因子を確かめやすいように,主に $u,v'$ を書いて計算する.

ここで1つ断りをいれておこう.本章では,不定積分の $C$ は任意定数を表し,現れる関数は考える区間で連続微分可能とする.対数を含む式は真数が正になる範囲で考える.以後この条件を本文では繰り返さないが,答案等に書く際は常に確かめること.

\begin{bookpoint}
\textbf{POINT}\par
微分すると簡単になる因子を $u$,原始関数が分かる因子を $v'$ と選び,\,$\int uv'\dd x=uv-\int u'v\dd x$ で残りの積分 $\int u'v\dd x$ へ移す.
\end{bookpoint}

\begin{example}{一次式と指数関数}
$\int x\e^x\dd x$ を求めよ.
\end{example}

\paragraphbf{方針}
$x$ を微分して $1$ にし,\,$\e^x$ を積分する.残る積分は $\int\e^x\dd x$ である.

\paragraphbf{解答}
$(x-1)\e^x+C$

\paragraphbf{解法}
$x$ は微分すると $1$ になり,\,$\e^x$ は積分しても $\e^x$ のままである.そこで
\[
 u=x,\qquad v'=\e^x,
 \qquad u'=1,\qquad v=\e^x
\]
と選ぶ.部分積分の公式より,
\[
\begin{aligned}
 \int x\e^x\dd x
 &=x\e^x-\int 1\cdot\e^x\dd x\\
 &=x\e^x-\e^x+C\\
 &=(x-1)\e^x+C
\end{aligned}
\]
を得る.答えを微分すると
\[
 \frac{\dd}{\dd x}\{(x-1)\e^x\}
 =\e^x+(x-1)\e^x
 =x\e^x
\]
となり,確かにもとの関数に戻る.

\paragraphbf{引いた項の意味}
章の初めで述べたように,積 $x\e^x$ を作って終わりにすると,その微分は $x\e^x+\e^x$ となる.例題の答えの $-\e^x$ は,この余分な $\e^x$ を取り除くために引いた項である.部分積分の公式の $-\int u'v\dd x$ は,一般にこの「余分な項」の原始関数を表しているのだ.

\paragraphbf{因子を逆に選んだ場合}
逆に $u=\e^x$,\,$v'=x$ と選ぶことも,公式上は正しい.その場合には
\[
 \int x\e^x\dd x
 =\frac{x^2}{2}\e^x-\frac12\int x^2\e^x\dd x
\]
となる.ところが,もとの一次式が二次式になり,残った積分はむしろ手間が増えている.公式を使えることと,その選択で計算が進むことは別である.右辺に何が残るかを一度書き下してから,因子の選び方を判断すればよい.
